% GENERATED FILE. Edit canonical lessons, then run npm run generate:books.
% canonical-source-hash: fnv1a64:fea3a0d732fcaaa9
% canonical-lessons: HI-C418-nazdik, HI-C418-nikat, HI-C418-par, HI-C418-bagal, HI-C418-pahucana

\chapter{Near and Across}
\label{ch:hi-near-and-across}
\hlchaptermodality{418}

\begin{chapteropening}
\textbf{By the end of this chapter:} \emph{I can say a place is close by, across the river or beside the bank, and take someone there.}

Complete the last lesson of chapter 418: I can say a place is close by, across the river or beside the bank, and take someone there.
\end{chapteropening}

\section[\texorpdfstring{\hi{नज़दीक} (nazdīk)}{नज़दीक (nazdīk)}]{\hi{नज़दीक} (nazdīk) — close by; near}

\label{lesson:HI-C418-nazdik}



\begin{quote}
Before the new one: say the Hindi for a route, then the Hindi for the way something faces.
\end{quote}

\subsection*{You'll want to know: \hi{नज़दीक}}
\textbf{\hi{नज़दीक}} — \emph{nazdīk} — \textquotedblleft{}close by; near\textquotedblright{}.

\begin{quote}
\hi{मेरा} \hi{घर} \hi{स्टेशन} \hi{के} \hi{नज़दीक} \hi{है।}
\end{quote}

\begin{quote}
\emph{merā ghar sṭeśan ke nazdīk hai.}
\end{quote}

\begin{quote}
My house is close to the station.
\end{quote}

\begin{cousinweb}
\textbf{\hi{नज़दीक}} (\emph{nazdīk}) is usually traced to Persian \emph{nazdīk}, \textquotedblleft{}near\textquotedblright{}; like \hi{पास} it follows \hi{के}: \hi{घर} \hi{के} \hi{नज़दीक}.
\end{cousinweb}

\begin{grammarlens}[title={What you've built}]
One more word for getting around.
\end{grammarlens}

\subsection*{Your turn}
\begin{itemize}
\raggedright
  \item \emph{Say it:} \emph{nazdīk}
  \item \emph{Say it:} \emph{nazdīk}, once more
  \item \emph{Say it:} \emph{merā ghar sṭeśan ke nazdīk hai.} — My house is close to the station.
  \item \emph{Recall:} say the Hindi for a route, then the Hindi for the way something faces, then say \emph{nazdīk} again
\end{itemize}

\begin{tcolorbox}[breakable,colback=teal!4,colframe=teal!35!black,title={Before you move on}]
What does \hi{नज़दीक} mean? (\textquotedblleft{}Close by; near\textquotedblright{}.) Say it once more.
\end{tcolorbox}
\section[\texorpdfstring{\hi{निकट} (nikaṭ)}{निकट (nikaṭ)}]{\hi{निकट} (nikaṭ) — near (formal)}

\label{lesson:HI-C418-nikat}



\begin{quote}
Before the new one: say the Hindi for the way something faces, then the Hindi for close by.
\end{quote}

\subsection*{You'll want to know: \hi{निकट}}
\textbf{\hi{निकट}} — \emph{nikaṭ} — \textquotedblleft{}near (formal)\textquotedblright{}.

\begin{quote}
\hi{होटल} \hi{स्टेशन} \hi{के} \hi{निकट} \hi{है।}
\end{quote}

\begin{quote}
\emph{hoṭal sṭeśan ke nikaṭ hai.}
\end{quote}

\begin{quote}
The hotel is near the station.
\end{quote}

\begin{cousinweb}
\textbf{\hi{निकट}} (\emph{nikaṭ}) is a tatsama, Sanskrit \emph{nikaṭa}, \textquotedblleft{}being at the side, near\textquotedblright{} (Monier-Williams); it is the formal partner of \hi{नज़दीक} and \hi{पास}.
\end{cousinweb}

\begin{grammarlens}[title={What you've built}]
One more word for getting around.
\end{grammarlens}

\subsection*{Your turn}
\begin{itemize}
\raggedright
  \item \emph{Say it:} \emph{nikaṭ}
  \item \emph{Say it:} \emph{nikaṭ}, once more
  \item \emph{Say it:} \emph{hoṭal sṭeśan ke nikaṭ hai.} — The hotel is near the station.
  \item \emph{Recall:} say the Hindi for the way something faces, then the Hindi for close by, then say \emph{nikaṭ} again
\end{itemize}

\begin{tcolorbox}[breakable,colback=teal!4,colframe=teal!35!black,title={Before you move on}]
What does \hi{निकट} mean? (\textquotedblleft{}Near (formal)\textquotedblright{}.) Say it once more.
\end{tcolorbox}
\section[\texorpdfstring{\hi{पार} (pār)}{पार (pār)}]{\hi{पार} (pār) — across; the far side}

\label{lesson:HI-C418-par}



\begin{quote}
Before the new one: say the Hindi for close by, then the Hindi for near (formal).
\end{quote}

\subsection*{You'll want to know: \hi{पार}}
\textbf{\hi{पार}} — \emph{pār} — \textquotedblleft{}across; the far side\textquotedblright{}.

\begin{quote}
\hi{हमने} \hi{नाव} \hi{से} \hi{नदी} \hi{पार} \hi{की।}
\end{quote}

\begin{quote}
\emph{hamne nāv se nadī pār kī.}
\end{quote}

\begin{quote}
We crossed the river by boat.
\end{quote}

\hi{पार} \hi{करना}, to cross; \hi{नदी} \hi{के} \hi{पार}, across the river; \hi{उस} \hi{पार}, on the far side.

\begin{cousinweb}
\textbf{\hi{पार}} (\emph{pār}) is a tadbhava of Sanskrit \emph{pāra}, \textquotedblleft{}the further bank, the far end\textquotedblright{} (CDIAL 8100), which lists Hindi \emph{pār}, \textquotedblleft{}on the other side of, across\textquotedblright{}.
\end{cousinweb}

\begin{grammarlens}[title={What you've built}]
One more word for getting around.
\end{grammarlens}

\subsection*{Your turn}
\begin{itemize}
\raggedright
  \item \emph{Say it:} \emph{pār}
  \item \emph{Say it:} \emph{pār}, once more
  \item \emph{Say it:} \emph{hamne nāv se nadī pār kī.} — We crossed the river by boat.
  \item \emph{Recall:} say the Hindi for close by, then the Hindi for near (formal), then say \emph{pār} again
\end{itemize}

\begin{tcolorbox}[breakable,colback=teal!4,colframe=teal!35!black,title={Before you move on}]
What does \hi{पार} mean? (\textquotedblleft{}Across; the far side\textquotedblright{}.) Say it once more.
\end{tcolorbox}
\section[\texorpdfstring{\hi{बग़ल} (baġal)}{बग़ल (baġal)}]{\hi{बग़ल} (baġal) — the side; beside (with \hi{में})}

\label{lesson:HI-C418-bagal}



\begin{quote}
Before the new one: say the Hindi for near (formal), then the Hindi for across.
\end{quote}

\subsection*{You'll want to know: \hi{बग़ल}}
\textbf{\hi{बग़ल}} — \emph{baġal} — \textquotedblleft{}the side; beside (with \hi{में})\textquotedblright{}.

\begin{quote}
\hi{बैंक} \hi{के} \hi{बग़ल} \hi{में} \hi{एक} \hi{दवाख़ाना} \hi{है।}
\end{quote}

\begin{quote}
\emph{baiṅk ke baġal meṁ ek davākhānā hai.}
\end{quote}

\begin{quote}
There is a pharmacy next to the bank.
\end{quote}

Feminine.

\begin{cousinweb}
\textbf{\hi{बग़ल}} (\emph{baġal}) is usually traced to Persian \emph{baġal}, \textquotedblleft{}the armpit, the side of the body\textquotedblright{}: what is in your armpit is right beside you.
\end{cousinweb}

\begin{grammarlens}[title={What you've built}]
One more word for getting around.
\end{grammarlens}

\subsection*{Your turn}
\begin{itemize}
\raggedright
  \item \emph{Say it:} \emph{baġal}
  \item \emph{Say it:} \emph{baġal}, once more
  \item \emph{Say it:} \emph{baiṅk ke baġal meṁ ek davākhānā hai.} — There is a pharmacy next to the bank.
  \item \emph{Recall:} say the Hindi for near (formal), then the Hindi for across, then say \emph{baġal} again
\end{itemize}

\begin{tcolorbox}[breakable,colback=teal!4,colframe=teal!35!black,title={Before you move on}]
What does \hi{बग़ल} mean? (\textquotedblleft{}The side; beside (with \hi{में})\textquotedblright{}.) Say it once more.
\end{tcolorbox}
\section[\texorpdfstring{\hi{पहुँचाना} (pahũcānā)}{पहुँचाना (pahũcānā)}]{\hi{पहुँचाना} (pahũcānā) — to take (someone) to; to deliver}

\label{lesson:HI-C418-pahucana}



\begin{quote}
Before the new one: say the Hindi for across, then the Hindi for beside.
\end{quote}

\subsection*{You'll want to know: \hi{पहुँचाना}}
\textbf{\hi{पहुँचाना}} — \emph{pahũcānā} — \textquotedblleft{}to take (someone) to; to deliver\textquotedblright{}.

\begin{quote}
\hi{ड्राइवर} \hi{ने} \hi{हमें} \hi{होटल} \hi{तक} \hi{पहुँचाया।}
\end{quote}

\begin{quote}
\emph{ḍrāivar ne hamẽ hoṭal tak pahũcāyā.}
\end{quote}

\begin{quote}
The driver took us all the way to the hotel.
\end{quote}

Past, with \hi{ने} (he or she …): \hi{उसने} \hi{चिट्ठी} \hi{घर} \hi{तक} \hi{पहुँचाई}, \emph{usne ciṭṭhī ghar tak pahũcāī}.

\begin{cousinweb}
\textbf{\hi{पहुँचाना}} (\emph{pahũcānā}) is the causative partner of \hi{पहुँचना}, \textquotedblleft{}to arrive\textquotedblright{}, which CDIAL files under Sanskrit \emph{prabhūta}, \textquotedblleft{}much, great\textquotedblright{} (CDIAL 8716): to make someone or something arrive.
\end{cousinweb}

\begin{grammarlens}[title={What you've built}]
One more word for getting around.
\end{grammarlens}

\subsection*{Your turn}
\begin{itemize}
\raggedright
  \item \emph{Say it:} \emph{pahũcānā}
  \item \emph{Say it:} \emph{pahũcānā}, once more
  \item \emph{Say it:} \emph{ḍrāivar ne hamẽ hoṭal tak pahũcāyā.} — The driver took us all the way to the hotel.
  \item \emph{Recall:} say the Hindi for across, then the Hindi for beside, then say \emph{pahũcānā} again
\end{itemize}

\begin{tcolorbox}[breakable,colback=teal!4,colframe=teal!35!black,title={Before you move on}]
What does \hi{पहुँचाना} mean? (\textquotedblleft{}To take (someone) to; to deliver\textquotedblright{}.) Say it once more.
\end{tcolorbox}
